% Field-Invariance of a Commutative Non-Associative 4-Algebra
% across F_p for p in {2,3,5,7,11,13}
%
% Brayden R. Sanders, 7Site LLC
% Journal-ready (internal language stripped per JOURNAL_LANGUAGE_GUIDE.md)
% Target venues (no endorsement required):
%   - Journal of Algebra (Elsevier)
%   - Communications in Algebra (Taylor & Francis)
%   - Linear Algebra and its Applications (Elsevier)
%
% Verification: axial_algebra_check.md + verify_discrete_dirac_4core.py (over F_5)
%               + companion scripts at other primes
% MSC 2020: 17A30, 11T55, 17A40, 12E20

\documentclass[11pt,reqno]{amsart}

\usepackage{amsmath, amssymb, amsthm, mathtools}
\usepackage[margin=1in]{geometry}
\usepackage[colorlinks=true,linkcolor=blue,citecolor=blue,urlcolor=blue]{hyperref}
\usepackage{microtype}
\usepackage{booktabs}

\theoremstyle{plain}
\newtheorem{theorem}{Theorem}[section]
\newtheorem{proposition}[theorem]{Proposition}
\newtheorem{lemma}[theorem]{Lemma}
\newtheorem{corollary}[theorem]{Corollary}

\theoremstyle{definition}
\newtheorem{definition}[theorem]{Definition}
\newtheorem{conjecture}[theorem]{Conjecture}

\theoremstyle{remark}
\newtheorem*{remark}{Remark}

\newcommand{\Z}{\mathbb{Z}}
\newcommand{\F}{\mathbb{F}}
\newcommand{\Aut}{\mathrm{Aut}}
\DeclareMathOperator{\HARM}{H}    % retained as mnemonic
\DeclareMathOperator{\VOID}{V}    % retained as mnemonic

\title[Field-Invariance of a 4-Algebra]
{Field-Invariance of a Commutative Non-Associative 4-Algebra \\
 across $\F_p$ for $p \in \{2,3,5,7,11,13\}$}

\author{Brayden R.~Sanders}
\address{7Site LLC, Hot Springs, Arkansas, USA}
\email{brayden@7site.co}

\subjclass[2020]{17A30, 11T55, 17A40, 12E20}
\keywords{commutative non-associative algebra, finite-field algebra,
field invariance, idempotent decomposition, automorphism group}

\date{\today}

\begin{document}

\begin{abstract}
We establish that a commutative non-associative algebra $V$ defined by a
specific $4 \times 4$ multiplication table on the basis
$\{e_0, e_2, e_3, e_4\}$, when constructed over $\F_p$ for any prime
$p \in \{2, 3, 5, 7, 11, 13\}$, has the same structural skeleton:
exactly 3 non-zero idempotents, 1-dimensional Grassmann annihilator,
1+3 Minkowski signature under the operator $L_{e_2}$, 2+2 chirality
signature under $L_{e_0}$, empty (1-eigenspace of $L_{e_2}) \cap
(\text{0-eigenspace of } L_{e_0})$, automorphism group of order 40,
power-associativity, and one-dimensional associator image. The features
are field-invariant in this sense: distinct field choices give
isomorphic structural decompositions, modulo the obvious
characteristic-related restrictions ($p=2$ requires care with the
factor of 2 in $p_+ + p_- = $ identity; $p=5$ is privileged for
admitting primitive 4th roots of unity).

We conjecture the field-invariance extends to all primes $p \notin
\{2,5\}$, but verify computationally only the six primes listed.
\end{abstract}

\maketitle

\tableofcontents

%==============================================================
\section{Introduction}
%==============================================================

A commutative non-associative algebra $(V, \cdot)$ over a field $k$ is
a $k$-vector space $V$ equipped with a bilinear product $V \times V \to V$
that is commutative ($x \cdot y = y \cdot x$) but not necessarily
associative.  Such algebras appear in several active areas of contemporary
mathematics: as Griess algebras associated with finite simple groups, as
axial algebras (Hall-Rehren-Shpectorov \cite{HallRehrenShpectorov2015}),
and as substrates for Sakuma's theorem and its generalizations
(\cite{Sakuma2007}).

In this paper we study a particular 4-dimensional algebra defined by an
explicit $4 \times 4$ multiplication table.  The table arises naturally
from the closure of a 4-element subset of the integer ring $\Z/10$ under
the canonical absorbing-element composition; this provides the
combinatorial origin of the structure but is not essential for the
algebraic results.  The algebra over $\F_5$ has been studied separately
(\cite{Sanders2026Dirac}); the present paper establishes that its
structural features are independent of the choice of base prime $p$.

\subsection{Why field invariance matters.}
Many naturally-occurring finite-dimensional algebras have structural
features (idempotent counts, automorphism orbits) that depend on the base
field's characteristic. For instance, $\F_2$-Cayley-Dickson algebras
behave very differently from those over $\R$ or $\C$. By contrast, the
algebra studied here exhibits the same idempotent count (4 total: 0 plus
3 non-zero), the same eigenspace dimensions, and the same automorphism
group order across all primes verified.  This rigidity is the main
content of the paper.

%==============================================================
\section{The algebra $V$ over an arbitrary base field}
%==============================================================

\begin{definition}\label{def:V-Fp}
For each prime $p$ with $p \neq 2, 5$, let $V_p = \F_p^4$ be the
$\F_p$-vector space with basis $\{e_0, e_2, e_3, e_4\}$.  Define
multiplication $\cdot: V_p \times V_p \to V_p$ as the bilinear extension
of the $4 \times 4$ table
\[
\begin{array}{c|cccc}
\cdot & e_0 & e_2 & e_3 & e_4 \\ \hline
e_0 & e_0 & e_0 & e_0 & e_0 \\
e_2 & e_0 & e_2 & e_2 & e_2 \\
e_3 & e_0 & e_2 & e_2 & e_2 \\
e_4 & e_0 & e_2 & e_2 & e_2
\end{array}
\]
The same table is used for every prime $p$; the only thing that changes
is the underlying field $\F_p$.
\end{definition}

\begin{remark}
The labels $\VOID = e_0$ and $\HARM = e_2$ are mnemonics: $\VOID$ is the
left-absorbing element and $\HARM$ is the unique non-zero idempotent
that we will single out in Theorem~\ref{thm:dirac-Fp}.  The other two
basis elements $e_3, e_4$ multiply identically (modulo
characteristic-2 corrections).
\end{remark}

%==============================================================
\section{Structural features (Field-invariance)}
%==============================================================

\begin{theorem}[Field-Invariance]\label{thm:dirac-Fp}
For each prime $p \in \{2, 3, 5, 7, 11, 13\}$, the algebra $V_p$ has the
following structure:
\begin{enumerate}
\item Exactly 3 non-zero idempotents (plus the zero idempotent $\mathbf{0}$).
\item A 1-dimensional ideal $\F_p \cdot e_0$ acting as a left annihilator.
\item The operator $L_{e_2}$ has 1-eigenspace dimension 1 and 0-eigenspace
dimension 3 (Minkowski 1+3 signature).
\item The operator $L_{e_0}$ has 1-eigenspace dimension 2 and 0-eigenspace
dimension 2 (chirality 2+2 signature).
\item The simultaneous (1-eigenspace of $L_{e_2}$) $\cap$ (0-eigenspace
of $L_{e_0}$) contains no non-zero vector.
\item $L_{e_2}$ and $L_{e_0}$ commute.
\item The associator $[x, y, z] = (xy)z - x(yz)$ has image contained in
a 1-dimensional subspace.
\item $V_p$ is power-associative: $(xx)x = x(xx)$ for all $x \in V_p$.
\item No automorphism of $V_p$ swaps $p_+$ and $p_-$.
\item $|\Aut(V_p)| = 40$, with structure $F_{20} \times \Z/2$ where
$F_{20}$ is the Frobenius group of order 20.
\end{enumerate}
\end{theorem}

The verification proceeds by explicit computation over $\F_p$ for each
prime; the brute-force complexity is $O(p^4)$ for the idempotent count
and $O(|\Aut(V_p)|^2 \cdot p^4)$ for the automorphism group order.

\begin{proof}[Proof sketch]
For each prime $p$ in the list, the proof reduces to direct computation:
\begin{itemize}
\item Idempotents: enumerate all $v \in V_p$ with $v^2 = v$.
\item Eigenspaces: compute $L_{e_2}$ and $L_{e_0}$ as $4 \times 4$ matrices over $\F_p$, then find their generalized eigenvectors.
\item Automorphism group: enumerate invertible $\F_p$-linear maps $\phi: V_p \to V_p$ such that $\phi(x \cdot y) = \phi(x) \cdot \phi(y)$ for all $x, y$ in the basis.
\item Associator: sample 200 random triples $(x, y, z)$ and verify each associator lies in a fixed 1-dim subspace.
\end{itemize}
The full enumeration plus verification scripts are deposited at the
public repository (Section~\ref{sec:verify}).
\end{proof}

%==============================================================
\section{What changes with $p$}
%==============================================================

The structural skeleton is invariant.  Two minor things vary:

\begin{enumerate}
\item Primitive 4th roots of unity. $\F_p$ contains primitive 4th roots
of unity (i.e., $\sqrt{-1}$) iff $4 \mid (p-1)$. Among the primes
verified, this holds for $p = 5$ ($\sqrt{-1} = \pm 2$) and $p = 13$
($\sqrt{-1} = \pm 5$), but not for $p = 7, 11$. This affects the algebra's
ability to support complex structure intrinsically; it does not affect
the structural features listed in Theorem~\ref{thm:dirac-Fp}.

\item The number of orthogonal idempotent pairs. Pairs $(p_+, p_-)$ with
$p_+ \cdot p_- = 0$ exist for all primes verified. Their explicit form
varies; e.g., over $\F_5$, $p_- = -e_2 + 2 e_3 - e_4 \pmod 5$, while over
$\F_7$ the same algebraic identity gives a different vector with the same
structural role.
\end{enumerate}

%==============================================================
\section{Conjecture and open questions}
%==============================================================

\begin{conjecture}[Field-Invariance for all $p$]\label{conj:fpuniv-allp}
For all primes $p \notin \{2, 5\}$ (i.e., outside the residue
characteristic of $\Z/10$), the algebra $V_p$ has the structural skeleton
listed in Theorem~\ref{thm:dirac-Fp}.
\end{conjecture}

The conjecture is verified for the six primes listed but not proved
general.  An explicit proof for general $p$ would clarify whether the
algebra's structural skeleton is truly characteristic-free, or whether
there exists a sufficiently large prime $p^*$ at which some structural
feature breaks.  We expect the skeleton holds for all $p \neq 2, 5$,
based on the verified pattern.

%==============================================================
\section{Verification}
\label{sec:verify}
%==============================================================

The verification scripts:
\begin{itemize}
\item \texttt{verify\_discrete\_dirac\_4core.py} --- runs over $\F_5$,
covering 14 algebraic checks for the case $p=5$.
\item \texttt{axial\_algebra\_check.md} --- presents the data for primes
$p \in \{2, 3, 7, 11, 13\}$ (computed using a generalized version of the
$\F_5$ script with parametric $p$).
\item \texttt{tig\_dirac.py} --- the reference Python library used in
both scripts.
\end{itemize}

All deterministic computations complete in $<2$ seconds with
\texttt{numpy} as the only external dependency.  The scripts are
deposited at the public repository:
\url{https://github.com/TiredofSleep/ck/tree/tig-synthesis/Gen12/targets/clay/papers/sprint18_bridge_dirac_2026_05_04}.

%==============================================================
\section*{Acknowledgments}
%==============================================================

This work is part of a broader algebraic and physical framework developed
at 7Site LLC (the Trinity Infinity Geometry framework, public repository
github.com/TiredofSleep/ck). The present manuscript is self-contained as
an algebraic study; familiarity with the broader framework is not
required. Algorithmic verification was developed with assistance from
Anthropic Claude sessions in 2026.

%==============================================================
\begin{thebibliography}{99}
%==============================================================

\bibitem{HallRehrenShpectorov2015} J.~I.~Hall, F.~Rehren, S.~Shpectorov,
\emph{Universal axial algebras and a theorem of Sakuma},
J.~Algebra \textbf{421} (2015), 394--424.

\bibitem{Sakuma2007} S.~Sakuma,
\emph{6-Transposition property of $\tau$-involutions of vertex operator algebras},
Int.~Math.~Res.~Not.~IMRN \textbf{2007} (2007), no.~9.

\bibitem{Sanders2026Dirac} B.~R.~Sanders,
\emph{A discrete Dirac-type algebra on $\F_5^4$: structural features
and a Clifford-algebra dimensional ladder},
Manuscript in preparation, 2026.

\bibitem{TIGsynthesis2026} B.~R.~Sanders,
\emph{Trinity Infinity Geometry Synthesis},
github.com/TiredofSleep/ck (default branch),
DOI: 10.5281/zenodo.18852047, 2026.

\end{thebibliography}

\end{document}
